% [原文 p.187]
量 $L$ 即变分问题的 \emph{Lagrange 量}（Lagrangian），是定义在 $M^n$ 的切丛全空间 $T(M^n)$ 的点 $(x,v)$ 上的实值光滑函数 $L(x,v)$，并且总要求它满足正定性条件
\begin{equation}
  \frac{\partial^2L}{\partial v^i\partial v^j}\,\eta^i\eta^j>0
  \qquad\text{对一切非零向量 }\eta. \tag{2.36}
\end{equation}
几何中或经典不带电粒子力学中出现的标准 Lagrange 量如下：
a) $L=\frac12g_{ij}v^iv^j=\frac12|v|^2$，这是 Riemann 流形 $M^n$ 中粒子运动的 Lagrange 量；只要参数 $t$ 是自然的（即与弧长成比例），极值曲线就是 $M^n$ 的参数化测地线；
b) $L=\frac12g_{ij}v^iv^j-U(x)$，适用于 Riemann 流形 $M^n$ 中粒子在有势 $U(x)$ 的力场中的运动；
c) $L(x,v)=|v|=\sqrt{g_{ij}v^iv^j}$，给出长度泛函，极值曲线为测地线；
d) 比 c) 更一般地，\emph{Finsler 长度}型泛函，其 Lagrange 量满足
\begin{equation}
  L(x,\lambda v)=\lambda L(x,v),\qquad\lambda>0; \tag{2.37}
\end{equation}
这类泛函出现在 Banach 空间的子流形上（取诱导度量）；通常要求它们满足正定性条件
\begin{equation}
  l(\gamma)=\int_\gamma L(x,\dot x)\,dt>0. \tag{2.38}
\end{equation}
若 $L(x,-v)=L(x,v)$，则称该变分问题是 \emph{可逆的}；例如关于 Riemann 度量的长度泛函 c) 就是如此。

按 Maupertuis 原理（或 Fermat--Jacobi--Maupertuis 原理），使 Lagrange 型泛函（见上面 b)）取极值的轨道，恰是关于由粒子总能量 $E=\frac12|v|^2+U(x)$ 定义的新度量的测地线：
\begin{equation}
  \begin{aligned}
    \tilde g^E_{ij}&=2(E-U)g_{ij}(x),\\
    l^E(\gamma)&=\int_\gamma\sqrt{2(E-U)g_{ij}\dot x^i\dot x^j}\,dt,
    \qquad \dot x^j=\frac{dx^j}{dt}.
  \end{aligned} \tag{2.39}
\end{equation}
只要 $E>\max_{x\in M^n}U(x)$，这显然定义了一个通常的 Riemann 度量。Fermat「最短时间原理」给出的极值曲线，基本上就是关于形如
\[
  g_{ij}=c^{-1}(x)\delta_{ij}\qquad\text{在 }\R^3\text{ 中}
\]
的度量的测地线，其中 $c(x)$ 是光滑纯量函数（最初是光速）。（众所周知，Fermat 提出光线沿最短时间的路径传播的原理，用以解释光线从一种各向同性介质进入另一种时所观察到的偏折，并把这一点与光在当时未知的、在不同介质中的不同速度联系起来。）
% [原文 p.188]

一般地，人们考虑定义在 \emph{道路空间} $\Omega(x_0,x_1)$ 上的泛函，该空间由 $M^n$ 中连接 $x_0$ 与 $x_1$ 的一切分段光滑参数化道路组成；在周期问题中则考虑空间 $\Omega^+(M^n)$，它由一切分段光滑、有向、参数化的闭路组成，即映射 $S^1\to M^n$，$f(\varphi+2\pi)=f(\varphi)$。假定流形 $M^n$ 闭。正定性条件 (2.38) 保证上述泛函 a)、b)、c)、d) 都满足紧性要求（「Arzelà 原理」）：对每个常数 $c$，满足 $S\{\gamma\}\le c$ 的道路 $\gamma$ 的子空间都相对紧（在 $\Omega(x_0,x_1)$ 或 $\Omega^+$ 的紧开拓扑下）。

关于紧开拓扑（或由 $M^n$ 上 Riemann 度量确定的道路间的自然距离），空间 $\Omega(x_0,x_1)$ 与 $\Omega^+$ 都是「无穷维」流形。由 Fermat--Jacobi--Maupertuis 原理，若只考虑闭流形 $M^n$ 上关于 Riemann 度量或 Finsler 型的长度泛函，上述例子 a)、b)、c)、d) 就基本上都被涵盖。再加上条件 (2.36)，就可以对 $\Omega(x_0,x_1)$（或 $\Omega^+$）上的这类泛函 $S\{\gamma\}$ 建立与前面所述流形上光滑函数的 Morse 理论类似的理论：用 $S$ 关于道路的「变分导数」代替梯度，于是临界道路就是变分导数 $\delta S/\delta\gamma$ 为零的道路；可以证明 $\gamma\in\Omega(x_0,x_1)$ 为临界道路当且仅当它是按自然参数参数化的光滑测地线。用泛函 $S$ 的「第二变分」$\delta^2S$，可把 Hessian $d^2S$ 定义为 $M^n$ 中沿 $\gamma$（$\gamma\in\Omega(x_0,x_1)$）的光滑向量场对（即 $T_\gamma(\Omega)$ 的元素）上的某个双线性形式；于是临界道路 $\gamma_0$ 的 \emph{指标} $i(\gamma_0)$ 定义为 $T_{\gamma_0}(\Omega)$ 的使 $d^2S$（作为 $T_{\gamma_0}(\Omega)$ 上的二次型）负定的子空间的最大维数。Morse 在 20 世纪 20 年代证明：临界道路 $\gamma_0\in\Omega(x_0,x_1)$ 的指标 $i(\gamma_0)$ 等于曲线 $\gamma_0$ 上与 $x_0$「共轭」的点 $p$ 的个数（计重数并除去 $x_1$）；所谓共轭，是指沿 $\gamma_0$ 从 $x_0$ 到 $p$ 的弧上存在非零向量场，它在 $x_0$ 与 $p$ 处为零且满足 Jacobi 线性微分方程（即由第二变分 $\delta^2S$ 给出的泛函的极值曲线的 Euler--Lagrange 方程）。在周期情形，即对 $\Omega^+(M^n)$，临界闭路上任一点处的共轭点个数只给出指标的下界估计。（这一理论的严格阐述特别需要 Hilbert 关于连接 $x_0$ 与 $x_1$ 的测地线存在性的定理。）与流形上函数的 Morse 理论一样，利用由「梯度」确定的轨道（以及用分段测地弧作有限维逼近），可以证明空间 $\Omega(x_0,x_1)$ 同伦等价于一个 $CW$-复形，其 $i$ 维胞腔与从 $x_0$ 到 $x_1$ 的指标为 $i$ 的测地线一一对应（对 $\Omega^+$ 类似）。由此可知，在一般的非退化情形（在上述假设下，包括 Arzelà 原理），此处也有 Morse 不等式
\begin{equation}
\begin{gathered}
  m_i(S)\ge b_i(\Omega(x_0,x_1)),\\
  m_i(S)\ge b_i(\Omega^+),
\end{gathered}\tag{2.40}
\end{equation}
% [原文 p.189]
其中 $m_i$ 是指标为 $i$ 的极值曲线的数目。

在某些情形下，第一个不等式可用来推断从 $x_0$ 到 $x_1$ 存在不同的测地线。但在周期问题中，多个闭测地线可能只是假象，例如它们可能全是一条简单闭测地线的重数（即沿该测地线绕不同圈数而得的闭路）。当 $M^n=S^n$ 且 $n\ge3$ 时，已知周期问题中至少存在两条几何不同的极值曲线（A. I. Fet，20 世纪 60 年代中期）。对任何闭流形 $M^n$，只要 Betti 数 $b_i(\Omega(x_0,x_1))$ 随 $i\to\infty$ 至少线性增长，就已知（Gromoll 与 Meyer，20 世纪 60 年代末）存在无穷多条几何不同的周期极值曲线。

在 Poincaré 关于 2-球面 $S^2$ 上三条不自交测地线的问题的背景下，Lyusternik 与 Shnirelman 把他们的不等式（见上）推广为 $m(S)\ge\operatorname{cat}(X/M^n)$。这里 $X$ 是 $M^n$ 上不自交道路空间的完备化，$M^n\subset X$ 是常值道路构成的子空间；$S$ 则是商空间 $X/M^n$ 上的任一可逆泛函。（这里同样不要求临界点非退化。）他们把问题转到无向道路空间，并证明 $X/S^2$（取 $M^n=S^2$）具有 $\RP^2$ 的同伦型，由此在 1920 年前后证明：对 $X/S^2$ 上任一可逆泛函 $S$，其临界点个数（不论退化与否）满足
\[
  m(S)\ge3,
\]
从而解决了 Poincaré 问题。\footnote{关于 Lyusternik--Shnirelman 理论的完整而严格的阐述最近已经发表；见 I. A. Taimanov 在 \emph{Russian Surveys}，1992 年，第 2 卷上的综述。}最后我们再次指出，这一理论只适用于可逆变分问题。

整体变分法的一个重要篇章，是研究由 Lie 群或对称空间上 Cartan--Killing 型度量确定的长度泛函与测地线（主要归功于 Bott，20 世纪 50 年代末）。例如在群 $SU_{2n}$ 上，赋以由 Lie 代数上的 Killing 形式 $\langle X,Y\rangle=\Re\Tr(X\bar Y^T)$ 确定的不变度量，从 $x_0=I$ 到 $x_1=-I$（$I$ 为单位矩阵）的极小测地线构成一个微分同胚于复 Grassmann 流形 $G^{\C}_{2n,n}\cong U_{2n}/(U_n\times U_n)$ 的子流形；作为作用量泛函 a) 的临界道路，极小测地线的指标为 0，而非极小测地线（从 $x_0$ 到 $x_1$）的指标 $\ge2n+2$。由此可知，极小测地线空间（$\cong G^{\C}_{2n,n}$）到 $\Omega(x_0,x_1)$ 的包含映射，诱导出相应同伦群之间的同构，只要这些同伦群的维数不超过（并包括）$2n$。

% [原文 p.190]
由主丛
\[
  V^{\C}_{2n,n}\to G^{\C}_{2n,n},\qquad\text{纤维为 }U_n
\]
（见 §1 式 (1.11)）的同伦正合列，推得
\begin{equation}
  \pi_j(G^{\C}_{2n,n})\cong\pi_{j-1}(U_n),\qquad 1\le j\le2n-2. \tag{2.41}
\end{equation}
把这些事实汇集起来（还须做些工作），就得到关于酉群稳定同伦群的著名结果，即 \emph{Bott 周期性}：
\begin{equation}
  \pi_j(U_n)\cong\pi_{j-2}(U_n),\qquad 2\le j\le2n-1. \tag{2.42}
\end{equation}
对正交群与辛群，论证要复杂一些。对正交群得到 8 重周期性：
\begin{equation}
  \pi_j(SO_n)\cong\pi_{j-8}(SO_n),\qquad 8\le j\le n-2. \tag{2.43}
\end{equation}
$O_n$ 与 $U_n$ 的前七个稳定同伦群由下表给出：
\begin{center}
\begin{tabular}{lcccccccccc}
\toprule
$j$ & 0 & 1 & 2 & 3 & 4 & 5 & 6 & 7 & 8 & $\cdots$\\
\midrule
$\pi_j(O_n)$ & $\Z/2$ & $\Z/2$ & 0 & $\Z$ & 0 & 0 & 0 & $\Z$ & $\Z/2$ & $\cdots$\\
$\pi_j(U_n)$ & 0 & $\Z$ & 0 & $\Z$ & 0 & $\Z$ & 0 & $\Z$ & 0 & $\cdots$\\
\bottomrule
\end{tabular}
\end{center}
对辛群 $Sp_n$ 有下述同构：
\begin{equation}
\begin{aligned}
  \pi_j(O_n)&=\pi_{j-4}(Sp_n),\qquad j\le n-2,\\
  \pi_j(Sp_n)&=\pi_{j-4}(O_n),\qquad j\le n-2.
\end{aligned}\tag{2.44}
\end{equation}
事实上，Bott 周期性算子对任一 $CW$-复形 $K$ 都有定义，即
\begin{equation}
\begin{aligned}
  {[\Sigma^8K,O_n]}&=[K,O_n]\qquad\text{若 }\dim K<n-10,\\
  {[\Sigma^2K,U_n]}&=[K,U_n]\qquad\text{若 }\dim K<2n-4,
\end{aligned}\tag{2.45}
\end{equation}
其中 $\Sigma$ 表示悬挂（suspension），$[X,Y]$ 表示映射 $X\to Y$ 的同伦类的集合。（通常把 (2.45) 中的下标 $n$ 略去，理解为它充分大。）

% [原文 p.191]
取正向极限
\[
  O=\lim_{n\to\infty}O_n,\qquad U=\lim_{n\to\infty}U_n,\qquad Sp=\lim_{n\to\infty}Sp_n,
\]
由 Bott 周期性推得
\[
  \Omega^4O=Sp,\qquad \Omega^4Sp=O,\qquad \Omega^2U=U,
\]
其中 $\Omega X=\Omega(X,x_0)$ 是 $X$ 上以 $x_0$ 为基点的闭路空间（基点常省略不写）。
% [Tutor Note: 这里的 $\Omega X$ 是闭路空间，与 p.188 起用于道路空间的 $\Omega(x_0,x_1)$、$\Omega^+$ 记法相近但含义不同：后者是无穷维流形，前者是拓扑空间的同伦不变量，阅读时勿混淆。]

\begin{remark}{}
不深入细节，我们提及一个事实（由好几位作者建立）：球面 $S^n$ 对 $n\ne1,3,7$ 的不可平行化性，可以由 Bott 周期性推出，而不必用 Adams 当初为此目的所用的结果。于是对一切 $n\ne1,3,7$，切丛 $\tau(S^n)$ 非平凡；而且当 $n$ 为奇数时，它（通过 $SO_n$-丛的分类——见第 3 章 §6）表示 $\pi_{n-1}(SO_n)$ 中一个 2 阶元素。注意对酉群 $U_n$，第一个非稳定同伦群是
\[
  \pi_{2n}(U_n)=\Z_{n!}.
\]
另一方面，Adams 关于不存在 Hopf 不变量为 1 的 $\pi_{2n-1}(S^n)$（$n\ge8$）元素的结果，似乎不能这样由 Bott 周期性推出（见第 3 章 §6）。\hfill$\square$
\end{remark}

下面转向 Lie 群、更一般地「H 空间」的上同调代数。\emph{H 空间}（H-space）是带连续乘法 $x\cdot y=\psi(x,y)\in X$（即 $\psi:X\times X\to X$）的拓扑空间 $X$，它关于该乘法有「同伦单位元」，即存在元素 $1\in X$，使定义乘以 1 的映射
\[
  \psi(1,x):X\to X,\qquad \psi(x,1):X\to X
\]
都同伦于恒等映射 $1_X$。注意 H 空间不必真是群，却往往「像」群：例如闭路空间 $\Omega(X,x_0)$ 的乘法只满足同伦意义下的结合律。闭路空间连同 Lie 群，构成 H 空间最重要的例子。例如由 Bott 周期性可知，空间 $BU=\lim_{n\to\infty}BU_n$ 具有闭路空间 $\Omega(U,1)$ 的同伦型：$BU\sim\Omega(U,1)$；对 $BO$ 与 $BSp$ 类似。因此就同调或同伦而言，它们可视为 H 空间。于是这些空间以及 $O$、$U$、$Sp$ 都是同伦结合且同伦交换的，尽管对有限的 $n$，即对 Lie 群 $O_n$、$U_n$、$Sp_n$，情形并非如此。

现在考虑任意 H 空间 $X$ 的上同调代数，系数取任意域 $k$。由自然的 $k$-代数同构
\[
  H^m(X\times X;k)\cong\sum_{j+q=m}H^j(X;k)\otimes H^q(X;k),
\]
以及 $X$ 上的 H 空间乘法 $\psi$，得到 $k$-代数同态
\begin{equation}
  \psi^*:H^*(X;k)\to H^*(X;k)\otimes H^*(X;k). \tag{2.46}
\end{equation}
% [原文 p.192]
由同伦单位元的存在推知，对任意 $z\in H^*(X;k)$，可把 $\psi^*(z)$ 写成
\begin{equation}
  \psi^*(z)=1\otimes z+z\otimes1+\sum_j z'_j\otimes z''_j, \tag{2.47}
\end{equation}
其中 $\dim z'_j,\dim z''_j>0$。条件 (2.46)、(2.47) 决定了 $H^*(X;k)$ 上的 Hopf 代数结构。形如 (2.47) 的代数同态的存在，对代数 $H^*(X;k)$ 的结构是很强的限制；按 Hopf 定理，特征 0 的域（例如 $\Q$、$\R$、$\C$）上分次、反称、结合的代数
\[
  H=\sum_{k\ge0}H^k,\qquad\dim H^k<\infty,
\]
若存在形如 (2.47) 的同态 $H\to H\otimes H$，则它是自由反交换分次代数，即除反交换性 $zw=(-1)^{ab}wz$（其中 $a$、$b$ 是 $z$、$w$ 的维数）外不满足任何关系。因此代数 $H^*(X;k)$ 有一组外生成元 $e_{j_\alpha}\in H^{2j+1}(X;k)$ 与一组多项式生成元 $u_{j_\beta}\in H^{2j}(X;k)$，它们共同生成 $H^*(X;k)$ 而无其他关系。由于自由多项式生成元 $u_{j_\beta}$ 满足 $u_{j_\beta}^n\ne0$ 对一切 $n$ 成立，立即可知：对有限维 H 空间 $X$（特别是 Lie 群），代数 $H^*(X;k)$ 是有限生成的自由外代数。

对特征 $p>0$ 的域 $k$，可以证明 $H^*(X;k)$ 有一组乘法生成元，用它们表示时只出现形如 $u^{p^n}=0$ 的关系。而且在若干情形已知：虽然 $u^{p^n}=0$，但在「对偶代数」中该维数处有非零元素。这里 H 空间 $X$ 的 Hopf 代数 $H^*(X;k)$ 的\emph{对偶（Pontryagin 对偶）代数}，是由同调群（向量空间）的直和 $H_*(X;k)$ 经同态
\[
  \psi_*:H_*(X;k)\otimes H_*(X;k)\to H_*(X;k)
\]
在其上定义乘法而得到的。

这些结果是靠对有限域上 H 空间上同调的细致考察得到的，这一方向由 Browder 在 20 世纪 60 年代初开创。

特别地，Browder 对有限维 H 空间的同调证明了 Poincaré 对偶律。然而有相当长一段时间不知道有限维 H 空间的非平凡例子，直到 20 世纪 70 年代初，Mislin、Hilton 等人利用 Quillen 与 Sullivan 约 1970 年为「Adams 猜想」发明的「同伦型范畴的 $p$-局部化技术」，才找到这类例子。（这一技术背后的想法其实可追溯到 Serre 在 20 世纪 50 年代前半期的工作。）这些新的 H 空间类似于某些经典 Lie 群与 H 空间对各种素数 $p$ 的同伦论 $p$-局部化。

% [原文 p.193]
对 $H$ 空间 $X=\Omega(Y,y_0)$，应当指出：在特征 0 的域 $k$ 上，Pontryagin 环 $H_*(X;k)$ 同构于同伦群
\[
  \pi_*(X)\otimes\Q,\qquad \pi_i(X)=\pi_{i+1}(Y)
\]
的 Lie--Whitehead 超代数的包络代数；分次方式与 $X$ 中相同，而元素的 Whitehead 乘法（Lie 超代数结构）$[x,y]$ 则与 $Y$ 中相同。为构造包络代数，取同伦群（张量 $\Q$）的某个分次基的全部元素 $x_j$，添入单位元 1，并考虑变量 $x$ 的一切非交换 $k$-多项式构成的代数，服从关系
\[
  x_ix_j-(-1)^{\dim x_i\dim x_j}x_jx_i=[x_i,x_j].
\]
这个结果是 Serre 与 Cartan 的定理的推论。

\begin{example}{}
对球面 $S^n$ 的闭路空间有
\begin{equation}
\begin{aligned}
  H^j(\Omega S^n;\Z)&=0,\qquad j\ne k(n-1),\\
  H^j(\Omega S^n;\Z)&\cong\Z,\qquad j=k(n-1),\ k\ge0.
\end{aligned}\tag{2.48}
\end{equation}
以 $v_k$ 记群 $H^{k(n-1)}(\Omega S^n;\Z)$ 的（适当选取的）生成元，则在 $H^*(\Omega S^n;\Z)$ 中有下述乘法关系：
\begin{equation}
\begin{aligned}
  1)\quad n=2k:\quad & v_1^2=0,\qquad
  v_{2p}v_{2q}=\binom{p}{p+q}v_{2(p+q)};\\
  2)\quad n=2k+1:\quad & v_pv_q=\binom{p}{p+q}v_{(p+q)}.
\end{aligned}\tag{2.49}
\end{equation}
同构 (2.48) 与关系 (2.49) 已被（Serre）证明可由相应纤维化的 Leray 上同调谱序列自动推出；不过它们最初是 Morse 与 Lyusternik 在 1930 年前后通过考察球面 $S^n$ 上的测地线得到的。
\end{example}
% [Tutor Note: 原书 (2.49) 的两处二项式系数均印作 $\binom{p}{p+q}$，而当 $p,q\ge1$ 时 $\binom{p}{p+q}=0$，这会使乘法恒为零，与 $\Omega S^n$ 上同调作为除幂代数的已知结构矛盾；按 2) 与 1) 的对称性及除幂代数关系 $v_pv_q=\frac{(p+q)!}{p!q!}v_{p+q}$，正确应为 $\binom{p+q}{p}$。此处照原书排印，批注说明。]

任意映射 $\psi:S^n\times S^n\to S^n$ 都可自然地（用 Hopf 的构造）延拓为映射
\[
  f_\psi:S^{2n+1}\to S^{n+1},
\]
这里借助下述标准嵌入把 $S^n$ 看作 $S^{n+1}$ 的子空间、把 $S^n\times S^n$ 看作 $S^{2n+1}$ 的子空间：$S^n$ 等同于 $S^{n+1}$ 的赤道：
\[
  S^{n+1}=D_1^{n+1}\cup D_2^{n+1},\qquad\text{其中 }S^n=D_1^{n+1}\cap D_2^{n+1};
\]
而 $S^n\times S^n$ 等同于 $S^{2n+1}$ 的下述子空间：
\[
  S^n\times S^n=(D_1^{n+1}\times S^n)\cap(S^n\times D_2^{n+1}),
\]
其中 $S^{2n+1}=(D_1^{n+1}\times S^n)\cup(S^n\times D_2^{n+1})$。

% [原文 p.194]
映射 $f_\psi$ 的「Hopf 不变量」$h(f,\psi)$ 恰好等于两个映射的度之积 $m_1m_2$：
\[
  \psi(s,s_0):S^n\times\{s_0\}\to S^n,\quad \deg\psi(s,s_0)=m_1,
\]
\[
  \psi(s_0,s):\{s_0\}\times S^n\to S^n,\quad \deg\psi(s_0,s)=m_2.
\]
若 $S^n$ 是带乘法 $\psi$ 的 H 空间，则显然两个度 $m_1,m_2=1$，于是存在 Hopf 不变量为 1 的映射 $f_\psi:S^{2n+1}\to S^{n+1}$。Hopf 证明这种映射只可能在 $n=2^r-1$ 时存在，而当 $n=1,3,7$ 时确实存在。Hopf 不变量为 1 的映射 $S^{2n+1}\to S^{n+1}$ 对 $n$ 不等于 1、3、7 时不存在，这一定理由 Adams 在 20 世纪 50 年代末证明。

这类定理——关于球面的不可平行化性，或上述 Adams 结果（它强得多）——作为推论给出：不存在维数不等于 2、4、8 的实有限维可除代数。（前面已指出，这也可由 Bott 周期性推出，正如 Milnor、Kervaire 与 Atiyah 所证明。）

20 世纪 80 年代初，在研究某些经典力学问题（例如重力场中刚体绕定点的运动、理想流体中自由运动的刚体）以及当代数学物理问题（特别是对 Dirac 磁单极子作细致的拓扑分析）的过程中，出现了一类多值泛函（例如作为 Dirac 磁单极子的推广与类比），它们定义在变量个数很大的某种道路空间或场空间上（Novikov）。多值函数与多值泛函的拓扑、其临界点个数的估计，以及其水平集结构的研究（有限维情形），构成本书下面将要描述的经典整体变分法的非经典类比。

有限维或无穷维流形 $M$ 上梯度单值的多值泛函，由 $M$ 上的闭 1-形式 $\omega$（$d\omega=0$）给出。多值函数 $S$ 本身则由积分
\[
  S(x)=\int_{x_0}^x\omega,\qquad\text{其中 }x_0\text{ 是固定点}
\]
给出。$S(x)$ 的值一般依赖于连接 $x$ 与 $x_0$ 的路径 $\gamma$。当把 $S$ 提升到某个正则无穷叶覆叠空间
\begin{equation}
  p:\tilde M\to M,\qquad\text{其中 }p^*\omega=dS \tag{2.50}
\end{equation}
上时，多值函数 $S$ 就成为单值的。假定 $H_1(M;\Z)\ne0$，这个覆叠空间的单值性群同构于 $\Z\times\Z\times\cdots\times\Z$（$k$ 个因子），其中 $(k-1)$ 称为形式 $\omega$ 的「无理度」。单值性群 $\Z\times\Z\times\cdots\times\Z$ 的生成元 $T_1,\dots,T_k$ 在 $\tilde M$ 上的作用由数 $\kappa_1,\dots,\kappa_k$ 确定，它们由「周期」给出。
% [原文 p.195]
首先取 $H_1(M;\Z)$ 的某个基 $\gamma_1,\dots,\gamma_m$，然后数 $\kappa_j$ 由形式 $\omega$ 沿 $\gamma_j$ 的积分（周期）给出：
\begin{equation}
\begin{aligned}
  \int_{\gamma_j}\omega&=\kappa_j,\qquad j\le k,\\
  \int_{\gamma_j}\omega&=0,\qquad j>k.
\end{aligned}\tag{2.51}
\end{equation}
（数 $\kappa_1,\dots,\kappa_k$ 在 $\Z$ 上线性无关。）上述单值性群的作用则由
\begin{equation}
  S(T_ix)=S(x)+\kappa_i,\qquad i=1,\dots,k, \tag{2.52}
\end{equation}
% [Tutor Note: 原书把 $k-1$（而非 $k$）称为 $\omega$ 的「无理度」：$k$ 是单值性群的秩，$k=1$ 即无理度为 0 的简单情形，此时群为 $\Z$，才有 (2.53) 的整上同调类。阅读时不要把 $k$ 直接当作无理度。]
给出。

若 $\{V_\alpha\}$ 是流形 $M$ 的开覆盖，且 $\omega$ 在每片上是恰当的：$\omega|_{V_\alpha}=dS^{(\alpha)}$，其中 $S^{(\alpha)}(x)$ 是 $V_\alpha$ 上的某个单值函数，则在每个重叠区域 $V_\alpha\cap V_\beta$ 上差 $S^{(\alpha)}-S^{(\beta)}$ 局部常值，即在 $V_\alpha\cap V_\beta$ 的每个连通分支上为常数。这给出构造流形 $M$（或任意拓扑空间）上梯度单值的多值函数 $S$ 的典范方法：取 $M$ 的开覆盖 $\{V_\alpha\}$ 及相应的一族函数 $\{S^{(\alpha)}\}$，每个 $S^{(\alpha)}$ 在 $V_\alpha$ 上有定义且单值，并具有性质：$S^{(\alpha)}-S^{(\beta)}$ 在 $V_\alpha\cap V_\beta$ 上局部常值。注意，若 $M$ 上存在单值函数 $S$，使 $S-S^{(\alpha)}$ 在每片 $V_\alpha$ 上为常数，则族 $\{S^{(\alpha)}\}$ 本质上就定义了这一单值函数。

若 $S$ 是 $M$ 上梯度单值的多值函数，则水平集 $S=\mathrm{const}$ 定义良好，并确定 $M$ 的一个「叶状结构」（可能有奇点），其「叶」一般可以非紧，即使 $M$ 紧。这类叶状结构具有相当大的拓扑意义，对它们的研究只是最近才开始。

与流形上单值光滑函数的情形一样，对多值函数 $S$（如上）人们关心临界点 $x_j$：$dS|_{x_j}=0$。无论在无穷维还是有限维情形，我们都只考虑这样的 1-形式 $\omega$：其临界点全是孤立、非退化的，且具有有限 Morse 指标（本质上就是二次型 $d^2S|_{x_j}$ 在 $x_j$ 处切空间上、关于适当局部坐标的负平方项的个数）。若单值性群为 $\Z$，生成元为 $\gamma_1$（换言之，1-形式 $\omega$ 的无理度为零），且其唯一的周期
\[
  \kappa_1=\int_{\gamma_1}\omega
\]
是整数（见上），则 $\omega$ 表示一个整上同调类
\begin{equation}
  [\omega]\in H^1(M;\Z). \tag{2.53}
\end{equation}
这确定了一个单值映射
\begin{equation}
  \psi=\exp\{2\pi iS\}:M\to S^1. \tag{2.54}
\end{equation}
% [原文 p.196]
在这种情形，称形式 $\omega$ 是\emph{量子化的}（quantized）。

\para{例 1.} 设 $M=M^n$ 是有限维闭流形，$\omega$ 是 $M$ 上的闭 1-形式。如何估计 $\omega$ 的临界点个数？上面提到的量子化形式的经典情形（见 (2.53)、(2.54)）特别简单：此时 $M^n$ 必须具有到 $S^1$ 上的纤维化的同伦型。关于无临界点的映射 $\psi:M\to S^1$ 的存在性，有各种拓扑判据，由 Browder、Levine、Hsiang 与 Farrell 在 20 世纪 60 年代得到（见下面 §5）。若 $M^n$ 不是这种类型，就会有临界点，其估计见下。

\para{例 2.} 设 $M$ 是道路空间 $\Omega(N^k,x_0,x_1)$，由 $N^k$（任意闭流形）中从 $x_0$ 到 $x_1$ 的道路组成；或设为有向闭路 $\gamma:S^1\to N^k$ 的空间 $\Omega^+(N^k)$。注意空间 $\Omega^+(N^k)$ 可视为 $N^k$ 上的纤维丛，投影映射为
\[
  p:\Omega^+(N^k)\to N^k,\qquad p(\gamma)=\gamma(0),
\]
它在点 $x\in N^k$ 上的纤维是闭路空间 $\Omega(N^k,x)$，由起止于 $x$ 的闭路组成。还有显然的截面 $\varphi:N^k\to\Omega^+(N^k)$，$p\varphi=1$，其中 $\varphi(x)$ 是 $x$ 处的常值闭路。若 $N^k$ 赋有 Riemann 度量，则 $N^k$ 上任一闭 2-形式 $\Omega$（对 $k\le3$ 可称为「广义磁场」，对 $k=4$ 称为「电磁场」）给出泛函
\begin{equation}
  S(\gamma)=\int_\gamma|\dot\gamma|^2dt+e\int_\gamma d^{-1}(\Omega) \tag{2.55}
\end{equation}
（「作用量泛函」），以及一个与之有相同极值曲线的泛函
\begin{equation}
  l^E(\gamma)=\int_\gamma E^{\frac12}|\dot\gamma|dt+e\int_\gamma d^{-1}(\Omega),
  \qquad E=\mathrm{const}. \tag{2.56}
\end{equation}
（「Maupertuis 泛函」）。这里 $E$ 是「能量」，$e$ 是「电荷」，$\gamma$ 取自 $\Omega(N^k,x_0,x_1)$ 或 $\Omega^+(N^k)$。只要 $N^k$ 单连通且形式 $\Omega$ 的上同调类非零：
\[
  0\ne[\Omega]\in H^2(N^k;\R),
\]
公式 (2.55)、(2.56) 就在空间 $M=\Omega(N^k,x_0,x_1)$（或 $\Omega^+(N^k)$）上定义多值泛函。
% [Tutor Note: 这里 $d^{-1}(\Omega)$ 指 $\Omega$ 的任一原像形式；因 $\Omega$ 闭，不同原像之差为闭形式，其沿道路的积分一般依赖道路的同伦类——这正是泛函取多值的来源。只有 $[\Omega]=0$ 时才能整体选到单值原像，这也解释了为何要假定 $0\ne[\Omega]$。]（若流形 $N^k$ 非单连通，则可将它嵌入（适当实现的）Eilenberg--MacLane 复形 $K(\pi,1)$，其中 $\pi=\pi_1(N^k)$，使得 $K(\pi,1)\setminus N^k$ 的所有胞腔维数 $\ge3$。诱导同态
\[
  H^2(\pi_1(N^k);\R)=H^2(K(\pi,1);\R)\to H^2(N^k;\R)
\]

% [原文 p.197]
是单射，由此可知在零伦自由闭路的子空间 $\Omega^+_0(N^k)$ 上，泛函 (2.55) 与 (2.56) 是单值的。）

\para{例 3.} 考虑光滑纤维丛
\[
  p:E\to B,
\]
其底空间 $B$ 是光滑闭流形或带边流形 $W^q$。设给定全空间 $E$ 上的闭 $(q+1)$-形式 $\Omega$（$d\Omega=0$），以及定义在纤维丛的截面 $\psi:B\to E$（$p\psi=1$）上的某个单值泛函 $S_0\{\psi\}$。公式
\begin{equation}
  S=S_0\{\psi\}+\int_{(B,\psi)}d^{-1}(\Omega), \tag{2.57}
\end{equation}
属于 (2.55) 与 (2.56) 的一般类型，它在丛的截面空间 $C$ 上定义一个泛函（一般是多值的），并在边界 $\partial W^q$（若有边界）上取预先给定的值。最有趣的情形是 $E=S^q\times G$，$G$ 为 Lie 群。此时截面就是所谓的「手征场」$g(x)$，$x\in\R^q$，且当 $|x|\to\infty$ 时 $g(x)\to g_0$。注意在每个紧单 Lie 群 $G$ 上总可定义非平凡的双侧不变 3-形式或 5-形式（例如当 $G=SU_n$，$n\ge3$ 时），分别适用于 $q=2$ 或 $4$ 的情形。\hfill$\square$

下面更详细地解释公式 (2.57)（以及 (2.55) 与 (2.56)）如何定义多值泛函。设 $\Omega$ 是纤维化 $p:E\to B$ 的全空间 $E$ 上的闭形式（如上），而 $\{V_\alpha\}$ 是 $E$ 的开覆盖，具有下述两条性质：
a) 形式 $\Omega$ 在每片 $V_\alpha$ 上恰当：
\[
  \Omega|_{V_\alpha}=d\Phi_\alpha;
\]
b) 对纤维化 $p:E\to B$ 的任一光滑截面 $\psi$，都存在指标 $\alpha$，使 $\psi(B)\subset V_\alpha$。
设 $C$ 是纤维化 $p:E\to B$ 的截面空间。形式 $\Phi_\alpha$ 在由含于 $V_\alpha$ 中的那些截面组成的区域 $W_\alpha\subset C$ 上确定「局部泛函」
\[
  S^{(\alpha)}(\psi)=\int_{(B,\psi)}\Phi_\alpha.
\]
在重叠区域上有
\[
  S^{(\alpha)}(\psi)-S^{(\beta)}(\psi)=\int_{(B,\psi)}(\Phi_\alpha-\Phi_\beta).
\]
由于 $d(\Phi_\alpha-\Phi_\beta)=0$ 在 $V_\alpha\cap V_\beta$ 上成立，按 Stokes 定理，差 $S^{(\alpha)}-S^{(\beta)}$ 在 $W_\alpha\cap W_\beta$ 上局部常值，因此族 $\{S^{(\alpha)}\}$ 确实确定 $C$ 上一个定义良好的 1-形式 $\omega$。

% [原文 p.198]
若在此背景下形式 $\omega$ 是量子化的，即 $[\omega]\in H^1(C;\Z)$，则所得（单值）映射
\[
  \exp\{2\pi iS\}:C\to S^1
\]
在 $\psi$ 处的值称为截面 $\psi$ 的 \emph{Feynman 振幅}。形式 $\omega$ 量子化这一条件，在带多值作用量泛函的这类场的量子理论中是本质的。它代表量子场论中所谓的「耦合常数的拓扑量子化」步骤，由 Novikov（1981）、Deser--Jackiw--Templeton 在 Chern--Simons 的特殊情形（1982）、以及 Witten（1983）所表述。

在平凡纤维化 $E=S^q\times F\to S^q$ 的情形，形式 $\Omega$ 实际上定义在纤维 $F$ 上，截面空间 $C$ 化为映射 $S^q\to F$ 的空间，而形式 $\omega$ 沿定义域泛函空间 $C$ 中道路的积分，与形式 $\Omega$ 在 $F$ 上沿 Hurewicz 同态 $\pi_{q+1}(F)\to H_{q+1}(F)$ 下适当像的积分一致。

对一般的丛 $E\to B$（如上），把 $\Omega$ 在 $E$ 中闭链上的积分与 $\omega$ 沿 $C$ 中道路的围道积分联系起来，需要依赖所论丛与边界条件作个别的拓扑分析（通常并不很难）。最后我们指出，若泛函 (2.57)（或 (2.55) 或 (2.56)）确实是多值的（即非单值），则在覆叠空间 $\tilde C$（参看 (2.50)）上——其中 $p_1^*\omega=\delta S$（$p_1:\tilde C\to C$）——泛函 $S$ 取遍一切值 $(-\infty,\infty)$。

下面转向多值函数与多值泛函的临界点个数的已知拓扑估计。从最简单的有限维情形开始。设 $w$ 是 $M^n$ 上无理度为零的闭 1-形式。于是有 $M^n$ 的覆叠空间 $\tilde M$，其单值性群为 $\Z$：
\[
  p:\tilde M\to M^n,\qquad p^*w=dS,
\]
其中生成单值性变换 $T:\tilde M\to\tilde M$ 满足
\[
  S(Tx)=S(x)+\kappa,\qquad\kappa\ne0.
\]
把 $M^n$ 与 $\tilde M$ 适当地分解为 $CW$-复形，就可赋予 $\tilde M$ 的胞腔链复形以 $\Z[T,T^{-1}]$-模结构（见第 3 章 §5）。用 $K^+$ 记环 $\Z[T,T^{-1}]$ 仅沿 $T$ 的正幂方向的完备化，即形如
\[
  q=\sum_{j\ge-N(q)}a_jT^j,\qquad a_j\in\Z
\]
的形式 Laurent 级数环。经表示（包含）$\rho:\Z[T,T^{-1}]\to K^+$，考虑 $\tilde M$ 的链复形以环 $K^+$ 为系数的同调群，这些同调群 $H_j(\tilde M;K^+)$ 于是可看作 $K^+$-模。由于 $K^+$ 是主理想整环，群 $H_j(\tilde M;K^+)$（我们可以写作 $H_j(\tilde M;[\omega])$，以表明它们依赖 $[\omega]$）有定义良好的无挠秩（Betti 数的类比）以及 $K^+$ 上的挠秩，分别记作
\[
  b_j(M^n,[\omega])=\mathrm{rk}\,H_j(\tilde M;K^+),
\]
% [原文 p.199]
与 $q_j(M^n,[\omega])$（挠子模的秩）。于是 $b_j+q_j$ 是 $K^+$-模 $H_j(\tilde M;K^+)$ 的最少生成元个数。只要多值函数 $S$ 的临界点孤立且全非退化，就能建立指标为 $i$ 的临界点个数 $m_i(S)$（或 $m_i(w)$）的下述下界（Novikov 不等式，20 世纪 80 年代初）：
\begin{equation}
  m_i(\omega)\ge b_i(M^n,[\omega])+q_i(M^n,[\omega])+q_{i-1}(M^n,[\omega]). \tag{2.58}
\end{equation}
较近（20 世纪 80 年代中期），Farber 通过构造适当的映射 $M^n\to S^1$（其中 $\pi_1(M^n)\cong\Z$，$n\ge6$）证明：这些不等式一般是严格的。

若无理度 $k-1\ge1$，从而覆叠空间 $\tilde M$ 的单值性群 $\Z\times\cdots\times\Z$ 多于一个因子，则情形更复杂。相应的环 $K^+_\kappa$ 依赖于周期组 (2.51)：$\kappa=(\kappa_1:\kappa_2:\dots:\kappa_k)$，它定义为环
\[
  \Z[T_1,\dots,T_k,T_1^{-1},\dots,T_k^{-1}]
\]
的扩充：把一切形式 Laurent 级数 $q$ 收进来，其中 $T_1^{m_1}T_2^{m_2}\cdots T_k^{m_k}$ 的系数 $a_{m_1m_2\dots m_k}$ 为零，除非对某个正整数 $N(q)$ 有 $\sum\kappa_jm_j>-N(q)$；并且对每一对数 $A<B$，只有有限多个非零系数 $a_{m_1\dots m_k}$ 满足 $A<\sum m_j\kappa_j<B$。同调群 $H_j(\tilde M;K^+_\kappa)$ 应当仍给出各指标临界点个数的某种下界。这类完备化的一种变体是环 $(K^+_\kappa)^{\mathrm{Stable}}$，那里还要求条件 $\sum\kappa_jm_j>-N(q)$ 不只是在特定的点 $\kappa=(\kappa_1:\kappa_2:\dots:\kappa_k)$ 处成立，而是在它的某个小邻域内成立。然而由于环 $K^+_\kappa$ 不再是主理想整环，确切的图景仍不清楚。\footnote{Sirokav 证明：对一般点 $\kappa$，环 $(K^+_\kappa)^{\mathrm{Stable}}$ 是整环，因此在更一般的情形 Morse 不等式也成立。最近 Pazhitnov 与 Ranicki 沿这些方向证明了进一步的结果，并在这一情形发展了 Morse 理论的类比（「Morse--Novikov 理论」）。}最后我们强调不等式 (2.58) 及其对 $k\ge2$ 的可能推广背后的关键事实：1-形式 $w$ 不仅给出一个链复形，而且给出一个具有同伦不变同调群的 $K^+_\kappa$-模复形。

下面转向函数空间 $\Omega(N^n,x_0,x_1)$ 与 $\Omega^+(N^n)$ 上 (2.55) 或 (2.56) 型泛函的类似问题。如我们所见，这些泛函确定闭 1-形式 $w=\delta S$。然而可以看出，对空间 $\Omega(N^n,x_0,x_1)$，用于得到不等式 (2.58) 的同调群全是退化的，因此对「固定端点的变分问题」，即对 $\Omega(N^n,x_0,x_1)$，似乎没有 Morse 不等式的类比。
% [原文 p.200]
另一方面，对「周期问题」，即对 $\Omega^+(N^n)$ 上的泛函，情形不同。在带 Riemann 度量 $g_{ij}$ 的完备 Riemann 流形 $N^n$ 上，例如 (2.39) 型的泛函（但允许为多值）总有这样的性质：由常值闭路构成的单点子流形 $\psi(N^n)\subset\Omega^+(N^n)$ 是泛函的局部极小点集。在多值情形，完全原像
\[
  p^{-1}(\psi(N^n))=\bigcup_j N^n_j,\qquad j=(j_1,\dots,j_k),\quad j_s\in\Z,
\]
在相应的覆叠空间 $\tilde M$（投影 $p:\tilde M\to M$）中，总由 $\tilde M$ 上单值函数 $\hat l^E$（由拉回 $p^*w=\delta\hat l^E$ 定义）的局部极小点组成。包含映射
\begin{equation}
  N^n\to N^n_j,\qquad \hat l^E(N^n_0)=0,\qquad
  \hat l^E N^n_j=\sum_{m=1}^k\kappa_mj_m, \tag{2.59}
\end{equation}
作为到 $\tilde M$ 的映射全都同伦。设 $\kappa_1>0$，考虑映射（同伦）
\[
  \Phi:N^n\times I\to F,\qquad I=[a,b],
\]
其中
\[
  \Phi(N^n\times\{a\})=N^n_{(0,\dots,0)}\subset\tilde M,\qquad
  \Phi(N^n\times\{b\})=N^n_{(1,0,\dots,0)}\subset\tilde M.
\]
有
\begin{equation}
  0<\min_\Phi\Bigl[\max_{x\in N^n}\hat l^E(\Phi(N^n\times I))\Bigr]
  <\hat l^E(N^n_{(0,\dots,0)})=0. \tag{2.60}
\end{equation}
% [Tutor Note: 原书 (2.60) 右端作 $\hat l^E(N^n_{(0,\dots,0)})=0$，但左端已严格大于 0，此不等式不可能成立；按上下文（考虑从 0 层到 $\kappa_1$ 层的同伦，其间必有临界点），右端应为 $\hat l^E(N^n_{(1,0,\dots,0)})=\kappa_1$。又 (2.60) 前一行把同伦的靶空间印作 $F$，按上下文应为覆叠空间 $\tilde M$。两处均照原书排印。]
本质上，把映射 $\Phi$ 在形如 $z\times I$（其中 $z$ 取遍 $H_*(N^n)$ 中的闭链）的子集上、沿与泛函 $\hat l^E$ 的梯度相反的方向形变，就得到区域 $l^E>0$ 中非退化临界点个数的下述不等式（假定处于一般位置）：
\begin{equation}
  m_{i+1}(l^E)\ge\max_p b_i^{(p)}(N^n), \tag{2.61}
\end{equation}
其中 $p$ 取遍所有素数（Novikov）。于是这些不等式就是 Morse 不等式 (2.8) 在「周期问题」中的类比。值得注意的是，不等式 (2.61) 不涉及函数空间 $\Omega^+(N^n)$ 本身的同调。然而由于某些分析上的困难，这些不等式在更一般的情形尚未建立。\footnote{Novikov 与 Grinevich（1994）对平坦 2-环面与负曲率曲面上非零磁场的情形建立了这些不等式。Ginzburg 从 20 世纪 80 年代末起用另一种技术研究了固定周期的磁场中的周期轨道（而不像先前的方法那样要求固定能量）。特别地，他帮助纠正了 Novikov 与 Taimanov 在 80 年代中期论文中的某些错误。}

% [原文 p.201]
不等式 (2.61) 也适用于 (2.38) 型泛函，只要 $l$ 是 $\Omega^+(N^n)$ 上的单值泛函，但「Arzelà 原理」被破坏，即存在闭路 $\gamma\in\Omega^+(N^n)$ 使 $l(\gamma)<0$（于是该度量构成 Finsler 型度量的非正类比；这类度量在数学物理中确实出现）。在这种情形，还需要利用（类似于上面）形如
\[
  \bar\Phi:[a,b]\to\Omega^+(N^n)
\]
的线段，使 $\bar\Phi(a)$ 是单点（即常值）闭路，而 $\bar\Phi(b)=\gamma$，且 $l(\gamma)<0$。有 Taimanov 的下述新近定理：

\emph{若对 Finsler 型的单值泛函（见 (2.37)）存在闭路 $\gamma\in\Omega^+(N^n)$ 使 $l(\gamma)<0$，且 $l$ 在所有常值闭路上为 0，则存在截面}
\[
  \psi_1:N^n\to\Omega^+(N^n),\qquad p\psi_1=1,
\]
\emph{使得}
a) \emph{$l(\psi_1(N^n))<0$，且}
b) \emph{截面 $\psi_1$ 经由同伦 $\bar\Phi:N^n\times I\to\Omega^+(N^n)$ 同伦于常值闭路的截面 $\psi$。}

这个结果建立了或可称为「闭链向泛函负值区域转移的原理」。对取负值也取正值的这类单值泛函 $l$，不等式 (2.61) 可由这一结果推出。关于闭路空间（或非自交曲线空间）上局部正则多值泛函的极值点的 Morse--Novikov 理论的新结果，可在 I. A. Taimanov 的综述中找到，见 \emph{Russian Surveys}，1992 年，第 2 卷。\footnote{到 20 世纪 80 年代末，辛几何专家（Floer、Hofer、A. Salomon、D. McDuff 等）已对自由闭路空间上的多值泛函建立了 Morse 理论。这些泛函一般不具有局部正则性，即临界点的 Morse 指标可以是无穷大。（这些想法最初出现在 Rabinovich 关于哈密顿系统周期轨道的工作中，20 世纪 80 年代初。）Floer 理论有对 3 维球面上联络模空间的应用，其中多值泛函与著名的 Chern--Simon 泛函一致。这一理论的最新发展还涉及上面定义的 Novikov 环 $K^+$。}
% [Tutor Note: 这条脚注里的「A. Salomon」原书 p.201 就是这样排印的；同一位作者在本附录 p.281 与正文其他处印作 Salamon（其本人的拼法是 Salamon），原书自身不一致，此处照录。]
% [Tutor Note: 原书 p.201 脚注 5 把该泛函印作 "Chern-Simon functional"（Simon 无 s），通行的写法是 Chern--Simons；此处照录原书拼法，读者按 Chern--Simons 理解即可。]

多值泛函的一个有趣例子（属于 Polyakov 与 Vigman）与 $\pi_3(S^2)$ 元素的 Hopf 不变量的「Whitehead 公式」有关；我们回忆该公式的构造：设 $f:S^3\to S^2$ 是光滑映射，$\omega$ 是 $S^2$ 上满足
\begin{equation}
  \int_{S^2}\omega=1 \tag{2.62}
\end{equation}
的 2-形式。设 $v$ 是 $S^3$ 上满足 $dv=f^*(\omega)=\bar\omega$ 的 1-形式。则 $f$ 的 Hopf 不变量的 Whitehead 公式为
\begin{equation}
  H\{f,\omega\}=\int_{S^3}v\wedge\bar\omega, \tag{2.63}
\end{equation}
% [原文 p.202]
并具有以下性质：
a) 同伦不变性：$\delta H/\delta f=0$；以及
b) 刚性：$\delta H/\delta\omega=0$（只要 (2.62) 成立）。

现在我们在零伦光滑映射 $f:S^2\to S^2$ 的空间 $F$ 上定义一个多值泛函（如前所述，由闭 1-形式确定）：
\begin{equation}
  S\{f\}=\int_{D^3_+}v\wedge\omega, \tag{2.64}
\end{equation}
其中映射 $f$ 已从 $S^2=\partial D^3_+$ 延拓到整个圆盘 $D^3_+$。$F$ 上相应的 1-形式则由 $\delta S$ 给出。

从 Chen 与 Sullivan 在 20 世纪 70 年代末的工作开始，人们构造了推广 (2.63) 的各种同伦不变积分。这类积分最清晰且最一般的版本（由 Novikov 在 20 世纪 80 年代初定义），还涉及「刚性」概念，对研究相应的多值泛函及其量子化很重要。下面我们描述这个积分。

设 $A=\sum_{i\ge0}A^i$ 是特征 0 的域 $k$ 上的反交换微分代数，$A^0=k$，上边缘算子 $d:A^i\to A^{i+1}$ 满足 $H^1(A)=0$。对每个正整数 $q$，我们构造 $A$ 的一个自由微分扩张 $C_q(A)$，它关于下述性质是极小的：
\begin{equation}
  H^j(C_q(A))=0,\qquad j\le q, \tag{2.65}
\end{equation}
构造如下：令 $C^0_q=A$。设 $\{y_{j\alpha}\}$ 是上同调环在维数 $j\le q$ 中的极小乘法生成元组。我们引入相应的自由生成元 $v_{j-1,\alpha}$，并令
\[
  C^1_q=A[\dots,v_{j-1,\alpha},\dots],\qquad dv_{j-1,\alpha}=y_{j\alpha}.
\]
注意包含映射 $A=C^0_q\to C^1_q$ 在维数 $\le q$ 上诱导上同调群的零同态。迭代这一构造，得到包含映射序列
\[
  A=C^0_q\subset C^1_q\subset C^2_q\subset\cdots\subset C_q(A).
\]
现在把代数 $A$ 的第 $(q+1)$ 个\emph{同伦群}（关于 $k$）定义为向量空间 $H^{q+1}(C^q(A))$：
\begin{equation}
  \pi_{q+1}(A)\otimes k=H^{q+1}(C_q(A)). \tag{2.66}
\end{equation}
% [Tutor Note: (2.66) 前一行正文把 $C_q(A)$ 印作 $C^q(A)$（上标），而 (2.66) 中印作 $C_q(A)$（下标）；按本节记号 $C^0_q\subset C^1_q\subset\cdots$ 与 (2.65)，正确的对象是下标族 $C_q(A)$，正文的上标应为排印之误。此处正文照原书排印。]
若取 $A=\Lambda^*(M^n)$，即单连通流形 $M^n$ 上 $C^\infty$-形式构成的代数，则可证明 $\pi_{q+1}(A)\otimes\R\cong\pi_{q+1}(M)\otimes\R$。因此任一单连通流形的同伦群的无穷部分可以相当初等地由 $M^n$ 上形式构成的代数 $\Lambda^*(M^n)$ 得到（尽管上述同构的完整验证需要代数拓扑中更精细的 Cartan--Serre 工具）。

% CHUNK ch04e DONE
